\scanpage{103}
\margin{IV.8}

$\psi^*\psi_*(\mathcal A')$ to $\mathcal A'$ corresponding to the identity homomorphism on $\psi_*(\mathcal A')$. We denote this map by $\sigma_{\mathcal A'}:\psi^*\psi_*(\mathcal A')\to\mathcal A'$. If $\mathcal A$ is a sheaf of rings on $X$, there is a canonical ring homomorphism from $\mathcal A$ to $\psi_*\psi^*(\mathcal A)$ corresponding to the identity homomorphism on $\psi^*(\mathcal A)$. We denote this map by $\rho_{\mathcal A}:\to\psi_*\psi^*(\mathcal A)$.

Now suppose $\Psi:(Y,\mathcal A')\to(X,\mathcal A)$ is a map of ringed spaces; i.e. $\Psi=(\psi,\theta)$ where $\psi$ is a continuous map from $Y$ to $X$ and $\theta$ is a homomorphism of sheaves of rings $\theta:\mathcal A\to\psi_*(\mathcal A')$. Then for any $\mathcal A'$-module $\mathcal M'$, $\psi_*(\mathcal M')$ can be made into an $\mathcal A$-module by means of the operation of $\theta$. When $\psi_*(\mathcal M')$ is considered a module over $\mathcal A$ in this way, it is denoted $\Psi_*(\mathcal M')$. We have the following canonical functorial homomorphisms [8, 0; \S\ 4.2]:
\[
h_*:\Psi_*(\uline{\mathrm{Hom}}_{\mathcal A'}(\mathcal M',\mathcal N'))\longrightarrow\uline{\mathrm{Hom}}_{\mathcal A}(\Psi_*(\mathcal M'),\Psi_*(\mathcal N'))
\]
\[
t_*:\Psi_*(\mathcal M')\underset{\mathcal A}{\otimes}\Psi_*(\mathcal N')\longrightarrow\Psi_*(\mathcal M'\underset{\mathcal A'}{\otimes}\mathcal N').
\]

On the other hand we have the map $\psi^*(\theta):\psi^*(\mathcal A)\to\psi^*\psi_*(\mathcal A')$, and so we have $\sigma_{\mathcal A'}:\psi^*(\theta):\psi^*(\mathcal A)\to\mathcal A'$. $\mathcal A'$ is an algebra over $\psi^*(\mathcal A)$ by the means of this map, and if $\mathcal M$ is an $\mathcal A$-module, we define the $\mathcal A'$-module $\Psi^*(\mathcal M)$ to be $\psi^*(\mathcal M)\underset{\psi^*(\mathcal A)}{\otimes}\mathcal A'$.

There is a canonical functorial isomorphism [8, 0; 4.33.1]
\[
t^*:\Psi^*(\mathcal M)\underset{\mathcal A'}{\otimes}\Psi^*(\mathcal N)\quad\Psi^*(\mathcal M\underset{\mathcal A}{\otimes}\mathcal N).
\]

It can be shown that there is a natural isomorphism from $\mathrm{Hom}_{\mathcal A'}(\Psi^*(\mathcal M),\mathcal M')$ to $\mathrm{Hom}_{\mathcal A'}(\mathcal M,\Psi_*(\mathcal M'))$ for all $\mathcal A$-modules $\mathcal M$ and $\mathcal A'$-modules $\mathcal M'$. By means of this isomorphism we can define the $\mathcal A$-homomorphism $\rho_{\mathcal M}:\mathcal M\to\Psi_*\Psi^*(\mathcal M)$ and the $\mathcal A'$-homomorphism

\mspage{102}
